Put \(h=h_n\), \(H=H_{n,\Delta}\), and \(\rho=\rho_{n,\Delta}\).  By \(\mathrm{def:bin\text{-}aware\text{-}estimator}\),
\[
 H=\Delta\left\lceil\frac{h}{\Delta}\right\rceil,
 \qquad h=n^{-1/4}.
\]
The inequalities \(x\leq\lceil x\rceil<x+1\), together with \(n\geq256\), \(0<\Delta\leq\Delta_0\), and \(\Delta_0\leq1/4\), give
\[
 h\leq H<h+\Delta\leq2\rho,
 \qquad 0<H<\frac12.
\]

We first establish the upper bounds using the single law-independent estimator in \(\mathrm{def:bin\text{-}aware\text{-}estimator}\).  Fix \(P\in\mathcal P_L\) and \(t\in\{0,1\}\).  Apply \(\mathrm{lem:conditional\text{-}local\text{-}binomial\text{-}law}\), and abbreviate its quantities by \(q_t\), \(\nu_t\), \(N_t\), and \(S_t\).  The class membership \(P\in\mathcal P_L\) and the displayed range of \(H\) imply
\[
 q_t\geq c_0H^2>0,
 \qquad |\nu_t-\theta_t|\leq LH,
 \qquad N_t\sim\operatorname{Binomial}(n,q_t).
\]
For every integer \(k\geq1\) having positive count probability, the same lemma gives
\[
 \mathbb E_P\!\left[\frac{S_t}{k}\,\middle|\,N_t=k\right]=\nu_t,
 \qquad
 \operatorname{Var}_P\!\left(\frac{S_t}{k}\,\middle|\,N_t=k\right)\leq\frac1{4k}.
\]
The selected outcomes are supported on \([0,1]\), so \(0\leq S_t/k\leq1\); consequently, the clipping in the estimator leaves \(S_t/k\) unchanged when \(N_t=k>0\).

We next derive the needed random-denominator inequalities.  If \(N\sim\operatorname{Binomial}(n,q)\) with \(q>0\), finite expansion and \((k+1)^{-1}=\int_0^1x^k\,dx\) yield
\[
 \mathbb E\!\left[\frac1{N+1}\right]
 =\int_0^1(1-q+qx)^n\,dx
 =\frac{1-(1-q)^{n+1}}{(n+1)q}
 \leq\frac1{(n+1)q}.
\]
Since \(k^{-1}\leq2(k+1)^{-1}\) for \(k\geq1\), and since \((\mathbb E U)^2\leq\mathbb E[U^2]\) for the nonnegative finite-valued variable \(U=\mathbf1\{N>0\}/\sqrt N\),
\[
 \mathbb E\!\left[\frac{\mathbf1\{N>0\}}N\right]
 \leq\frac2{(n+1)q}
 \leq\frac2{nq},
 \qquad
 \mathbb E\!\left[\frac{\mathbf1\{N>0\}}{\sqrt N}\right]
 \leq\sqrt{\frac2{nq}}.
\]
Furthermore, \(\mathbb P(N=0)=(1-q)^n\leq e^{-nq}\).  The elementary inequalities \(x^{1/2}e^{-x}\leq1\) and \(xe^{-x}\leq1\), valid for \(x>0\), therefore show
\[
 \mathbb P(N=0)\leq(nq)^{-1/2},
 \qquad
 \mathbb P(N=0)\leq(nq)^{-1}.
\]

The conditional variance inequality and \((\mathbb E|V|)^2\leq\mathbb E[V^2]\) imply, on \(N_t=k>0\),
\[
 \mathbb E_P\!\left[\left|\frac{S_t}{k}-\nu_t\right|\,\middle|\,N_t=k\right]
 \leq\frac1{2\sqrt k},
 \qquad
 \mathbb E_P\!\left[\left(\frac{S_t}{k}-\nu_t\right)^2\,\middle|\,N_t=k\right]
 \leq\frac1{4k}.
\]
On \(N_t=0\), the fixed fallback is \(1/2\), and \(0\leq\nu_t\leq1\) gives
\[
 |1/2-\nu_t|\leq\frac12,
 \qquad (1/2-\nu_t)^2\leq\frac14.
\]
Averaging over \(N_t\) and applying the random-denominator bounds with \(q=q_t\) proves
\[
 \mathbb E_P|\widehat m_{t,n,\Delta}-\nu_t|
 \leq\left(\frac1{\sqrt2}+\frac12\right)\frac1{\sqrt{nq_t}},
\]
and
\[
 \mathbb E_P[(\widehat m_{t,n,\Delta}-\nu_t)^2]
 \leq\frac3{4nq_t}.
\]
Let \(A=1/\sqrt2+1/2\).  The inequalities \(q_t\geq c_0H^2\), \(|\nu_t-\theta_t|\leq LH\), and \((a+b)^2\leq2a^2+2b^2\) now give
\[
 \mathbb E_P|\widehat m_{t,n,\Delta}-\theta_t|
 \leq LH+\frac{A}{\sqrt{c_0n}\,H},
\]
\[
 \mathbb E_P[(\widehat m_{t,n,\Delta}-\theta_t)^2]
 \leq2L^2H^2+\frac3{2c_0nH^2}.
\]
Because \(H\geq h=n^{-1/4}\),
\[
 \frac1{\sqrt n\,H}\leq h\leq\rho,
 \qquad
 \frac1{nH^2}\leq h^2\leq\rho^2.
\]
Together with \(H<2\rho\), the identity \(\tau_P(b)=\theta_1-\theta_0\), the triangle inequality, and \((x-y)^2\leq2x^2+2y^2\), these estimates imply
\[
 \sup_{P\in\mathcal P_L}
 \mathbb E_P|\widehat\tau_{n,\Delta}-\tau_P(b)|
 \leq C_1\rho,
 \qquad
 C_1=4L+\frac{2A}{\sqrt{c_0}},
\]
and
\[
 \sup_{P\in\mathcal P_L}
 \mathbb E_P[(\widehat\tau_{n,\Delta}-\tau_P(b))^2]
 \leq C_2\rho^2,
 \qquad
 C_2=32L^2+\frac6{c_0}.
\]
The counts and sums are measurable functions of the released sample \(Z_\Delta^n\), the bandwidth is known, and the zero-count fallback is fixed.  Hence \(\widehat\tau_{n,\Delta}\in\mathcal T_{n,\Delta}\).  Taking the two infima in \(\mathrm{def:minimax\text{-}risks}\) proves the asserted upper bounds and their attainment, up to constants, by the displayed estimator.

We now separate the population and sampling meanings of the converse.  By \(\mathrm{prop:release\text{-}fiber\text{-}diameter}\),
\[
 2c_L\Delta\leq\omega_\Delta\leq4L\Delta.
\]
Thus, for every fixed \(\Delta>0\), the structural full-data functional \(\tau_P(b)\) is generally not point identified by the released one-unit law; the \(\Delta\) term is a population release-fiber ambiguity cost.  In contrast, \(h=n^{-1/4}\) is the local sampling-error scale.  A single same-class pair carrying both scales is supplied by \(\mathrm{thm:single\text{-}joint\text{-}pair}\).

Let \(P_0=P_{0,n,\Delta}\) and \(P_1=P_{1,n,\Delta}\) be that pair, let \(Q_j=\mathcal L_{P_j}(Z_\Delta^n)\), and put \(\vartheta_j=\tau_{P_j}(b)\).  With \(\kappa=c_L/4>0\) and \(\delta=|\vartheta_1-\vartheta_0|\), the theorem gives
\[
 P_0,P_1\in\mathcal P_L,
 \qquad
 \operatorname{TV}(Q_0,Q_1)\leq\frac14,
 \qquad
 \delta=\kappa(h+\Delta)\geq\kappa\rho.
\]
Fix any \(a\in\mathcal T_{n,\Delta}\).  Set \(\lambda=Q_0+Q_1\) and \(f_j=dQ_j/d\lambda\).  For every released sample \(z\),
\[
 |a(z)-\vartheta_0|+|a(z)-\vartheta_1|\geq\delta,
\]
and
\[
 (a(z)-\vartheta_0)^2+(a(z)-\vartheta_1)^2
 =2\left(a(z)-\frac{\vartheta_0+\vartheta_1}{2}\right)^2+\frac{\delta^2}{2}
 \geq\frac{\delta^2}{2}.
\]
Moreover, directly from the density formula for total variation,
\[
 \int\min(f_0,f_1)\,d\lambda
 =\frac12\int(f_0+f_1-|f_0-f_1|)\,d\lambda
 =1-\operatorname{TV}(Q_0,Q_1).
\]
Since \(f_j\geq\min(f_0,f_1)\), multiplication of the two pointwise loss inequalities by \(\min(f_0,f_1)\), followed by integration and the fact that a maximum is at least half the corresponding sum, yields
\[
 \max_{j\in\{0,1\}}\mathbb E_{Q_j}|a-\vartheta_j|
 \geq\frac\delta2\{1-\operatorname{TV}(Q_0,Q_1)\}
 \geq\frac{3\kappa}{8}\rho,
\]
and
\[
 \max_{j\in\{0,1\}}\mathbb E_{Q_j}(a-\vartheta_j)^2
 \geq\frac{\delta^2}{4}\{1-\operatorname{TV}(Q_0,Q_1)\}
 \geq\frac{3\kappa^2}{16}\rho^2.
\]
The supremum over \(\mathcal P_L\) dominates the maximum over \(P_0,P_1\).  Taking the infimum over \(a\in\mathcal T_{n,\Delta}\) proves both minimax lower bounds.  Finally, choose
\[
 c=\min\left\{\frac{3\kappa}{8},\frac{3\kappa^2}{16}\right\}>0,
 \qquad
 C=1+\max\{C_1,C_2,c\}.
\]
These constants satisfy \(c<C<\infty\) and depend only on the permitted class constants.  All inequalities are uniform over \(n\geq256\) and \(0<\Delta\leq\Delta_0\), which completes the proof.